% 原书第六篇6.5节：放射性214--226；印刷页257，PDF266。
% 底本书末答案印刷348--349，PDF357--358。
\radioactivityproblem{214}{257}
放射性核衰变常数为 $\lambda$。
求 (a) 由 $t=0$ 到 $t$ 衰变的概率；
(b) 核的平均寿命。
\exsolution
单核在 $[0,t]$ 内不衰变的生存概率
$S(t)=e^{-\lambda t}$，
因此
\[
 \boxed{P(0<T_d<t)=1-e^{-\lambda t}}.
\]
衰变时间密度
$f(t)=\lambda e^{-\lambda t}$，故
\[
 \boxed{\tau=\int_0^\infty t f(t)\,dt
 =\int_0^\infty e^{-\lambda t}dt
 =\frac1{\lambda}}.
\]
\excheck 概率极限 $P(0)=0$、
$P(\infty)=1$；均寿命不等于半衰期。

\radioactivityproblem{215}{257}
半衰期 $T_{1/2}=71.3\,\mathrm d$ 的
放射性钴原子核，一个月内有多少份额发生衰变？
\exsolution
衰变常数
$\lambda=\ln2/T_{1/2}$。
按一个月约30天，
剩余份额 $2^{-30/71.3}$，
所以
\[
 \boxed{f=1-2^{-30/71.3}\simeq0.253
 \approx\frac14}.
\]
\excheck “一个月”并非固定秒数，
采用30天与实际公历月的近似差小于本题精度。

\radioactivityproblem{216}{257}
质量 $m=1.0\,\mu\mathrm g$ 的
$\mathrm{^{24}Na}$，半衰期15小时，
一小时内发射多少粒 $\beta$ 粒子？
\exsolution
初始核数
$N_0=(m/M)N_A$，
$M\simeq24\,\mathrm{g/mol}$。
假定每次 $\beta$ 衰变给出一电子，
一小时内衰变数
\[
 \boxed{\Delta N=
 \frac{mN_A}{M}
 \left(1-2^{-1/15}\right)
 \simeq1.13\times10^{15}\approx1.2\times10^{15}}.
\]
\excheck 不可用 $A_0\,t$ 当作精确值，
因为一小时内活度本身也随时间减少。

\radioactivityproblem{217}{257}
从 $t=0$ 开始记录
$\mathrm{^{23}Mg}$ 的 $\beta$ 衰变。
$t_1=2.0\,\mathrm s$ 内计得 $N_1$ 粒；
到 $t_2=3t_1$ 时累计计数
$N_2=2.66N_1$。
求平均寿命。
\exsolution
$N(t)=N_0(1-e^{-\lambda t})$。
令 $x=e^{-\lambda t_1}$，
由于 $t_2=3t_1$，
\[
 2.66=\frac{1-x^3}{1-x}=1+x+x^2.
\]
取 $0<x<1$ 的根
$x=[-1+\sqrt{1+4(2.66-1)}]/2\simeq0.882$。
故
\[
 \boxed{\tau=\frac1\lambda=
 -\frac{t_1}{\ln x}\simeq15.9\,\mathrm s
 \approx16\,\mathrm s}.
\]
\excheck 此题是\emph{累计计数}的比例，
不是瞬时活度的比例。

\radioactivityproblem{218}{257}
某放射性制剂的活度经过7.0天降低为原来的
$1/2.5$。求半衰期。
\exsolution
$A(t)/A(0)=e^{-\lambda t}=1/2.5$，
所以 $\lambda=\ln2.5/7.0\,\mathrm{d^{-1}}$。
由 $T_{1/2}=\ln2/\lambda$，
\[
 \boxed{T_{1/2}=
 \frac{7.0\,\mathrm d\,\ln2}{\ln2.5}
 \simeq5.30\,\mathrm d}.
\]
\excheck 活度与未衰变核数均随同一指数衰减。

\radioactivityproblem{219}{257}
某放射性核素初始活度为
$A_0=650\,\mathrm{min^{-1}}$。
经过\emph{半个半衰期}后活度是多少？
\exsolution
$A(t)=A_02^{-t/T_{1/2}}$。
当 $t=T_{1/2}/2$，
\[
 \boxed{A=650\,2^{-1/2}
 \simeq4.60\times10^2\,\mathrm{min^{-1}}}.
\]
\excheck 经过半个半衰期不是活度减半；
应乘 $2^{-1/2}$。

\radioactivityproblem{220}{257}
某 $\mathrm{^{55}Co}$ 制剂的活度每小时下降4.0\%，
产物不再有放射性。求衰变常数与平均寿命。
\exsolution
一小时后活度比
$A(1\,\mathrm h)/A_0=0.960$，
故
\[
 \boxed{\lambda=-\frac{\ln0.960}{3600\,\mathrm s}
 \simeq1.134\times10^{-5}\,\mathrm{s^{-1}}}.
\]
由 $\tau=1/\lambda$，
\[
 \boxed{\tau\simeq8.82\times10^4\,\mathrm s
 =24.5\,\mathrm h\simeq1.02\,\mathrm d}.
\]
\excheck 本版书末将均寿命写成“1.0 years”；
照印刷的\emph{每小时下降4\%}应得约1.0天，
而不可能是一年。题目参数与书末单位冲突，已记录。

\radioactivityproblem{221}{257}
质量1.0g的$\mathrm{^{238}U}$制剂
每秒发射 $1.24\times10^4$ 粒 $\alpha$ 粒子。
求半衰期及制剂活度。
\exsolution
一克纯 $\mathrm{^{238}U}$ 核数
$N=(1\,\mathrm g/238\,\mathrm{g/mol})N_A$。
若每次衰变发射一粒 $\alpha$，
$A=\lambda N$，所以
\[
 \boxed{T_{1/2}=\frac{N\ln2}{A}
 \simeq4.49\times10^9\,\mathrm{year}}.
\]
活度即给出的衰变率
\[
 \boxed{A=1.24\times10^4\,\mathrm{Bq}}.
\]
\excheck $\mathrm{Bq}=1\,\mathrm{s^{-1}}$；
不可把 $\alpha$ 计数与单位换算再乘一分钟。

\radioactivityproblem{222}{257}
古木中 $\mathrm{^{14}C}$ 比活度为新砍伐木材
比活度的 $3/5$，
已知 $\mathrm{^{14}C}$ 半衰期为5570年，
估算古木年代。
\exsolution
假定碳交换停止之后 $\mathrm{^{14}C}$ 只衰变，
\[
 \frac{A(t)}{A(0)}=e^{-\lambda t}=\frac35,\quad
 \lambda=\frac{\ln2}{5570\,\mathrm{yr}}.
\]
故
\[
 \boxed{t=\frac{5570}{\ln2}\ln\frac53
 \simeq4.11\times10^3\,\mathrm{yr}}.
\]
\excheck 需要假设最初与同环境新木的
$\mathrm{^{14}C}$ 比活度相同；
否则须进行年代校准。

\radioactivityproblem{223}{257}
某铀矿中 $\mathrm{^{238}U}$ 核数与
$\mathrm{^{206}Pb}$ 核数之比为
$\eta=2.8$。
假定矿中所有 $\mathrm{^{206}Pb}$
均来自该铀衰变链，
$T_{1/2}=4.5\times10^9$ 年，
估算铀矿年龄。
\exsolution
无初始铅且铀衰变链子代相对迅速达到稳定铅时，
$N_U=N_0e^{-\lambda t}$、
$N_{\rm Pb}=N_0-N_U$。
由 $N_U/N_{\rm Pb}=\eta$ 得
$e^{\lambda t}=1+1/\eta$，因此
\[
 \boxed{t=\frac{T_{1/2}}{\ln2}
 \ln\left(1+\frac1\eta\right)
 \simeq2.00\times10^9\,\mathrm{yr}}.
\]
\excheck 必须假定初始没有额外稳定铅，
否则不能仅凭测量的核数比唯一确定矿龄。

\radioactivityproblem{224}{257}
计算 $\mathrm{^{24}Na}$ 与 $\mathrm{^{235}U}$
的比活度，已知半衰期分别为15小时与
$7.1\times10^8$ 年。
\exsolution
一克纯放射性核素中
$N=N_A/M$，而 $A=\lambda N$。
故比活度
$a_s=(\ln2/T_{1/2})(N_A/M)$。
代入 $M_{\rm Na}=24\,\mathrm{g/mol}$、
$M_{\rm U}=235\,\mathrm{g/mol}$，
\[
 \boxed{a_s(\mathrm{^{24}Na})\simeq
 3.2\times10^{17}\,\mathrm{Bq/g}},
\quad
 \boxed{a_s(\mathrm{^{235}U})\simeq
 8.0\times10^{4}\,\mathrm{Bq/g}}.
\]
\excheck 比活度单位为单位质量每秒衰变数；
年代须先换成秒。

\radioactivityproblem{225}{257}
把活度 $A_0=2.0\times10^3\,\mathrm{Bq}$
的 $\mathrm{^{24}Na}$ 溶液注入血液。
$t=5.0$小时后，
每 $\mathrm{cm^3}$ 血液样本活度为
$A'=16\,\mathrm{min^{-1}\,cm^{-3}}$。
该核素半衰期15小时，求血液总体积。
\exsolution
假定核素在血液中均匀分布且无体外排泄，
$t$时剩余总活度为
$A_0\,2^{-t/T_{1/2}}$。
每立方厘米样本活度
$A'=16/60\,\mathrm{Bq/cm^3}$，
所以
\[
 \boxed{V=\frac{A_0}{A'}\,2^{-t/T_{1/2}}
 =\frac{2000}{16/60}\,2^{-5/15}
 \simeq5.95\times10^3\,\mathrm{cm^3}
 \approx6.0\,\mathrm L}.
\]
\excheck 实际医学使用须考虑组织分布、
排泄及剂量安全；这里纯属理想示踪守恒模型。

\radioactivityproblem{226}{257}
含放射性 $\mathrm{^{58}Co}$ 和非放射性
$\mathrm{^{59}Co}$ 的制剂，
整体比活度为
$2.2\times10^{12}\,\mathrm{Bq/g}$。
$\mathrm{^{58}Co}$ 半衰期为71.3天，
求放射性钴所占质量百分数。
\exsolution
若每克纯 $\mathrm{^{58}Co}$ 的比活度
$a_{58}=(\ln2/T_{1/2})(N_A/58)$，
混合物的整体比活度为 $a=f a_{58}$，
其中 $f$ 为放射性同位素质量份额。
因此
\[
 \boxed{f=\frac{a}{a_{58}}
 =\frac{aT_{1/2}\,58}{N_A\ln2}
 \simeq1.88\times10^{-3}},
 \quad \boxed{100f\simeq0.188\%\approx0.19\%}.
\]
\excheck 非放射性的 $\mathrm{^{59}Co}$
提供质量但不提供活度，所以直接按质量份额线性缩放。
